\scanpage{139}
% archive foliation: 138
\deleted{par la formule $z_2/z_1$, qu'on} pourrait penser $z=z_2/z_1$, \deleted{\unsure{et}} mais par $(24\deleted{\unsure{bis}})$ i.e. par
\[
z=-\unsure{\overline{(z_2/z_1)}|(z_2/z_1)^{-1}}
=-\unsure{\overline{z_2}/\overline{z_1}|z_2/z_1}\in\mathcal D.
\]
Si on voulait repérer \inserted{non} la structure de Torelli par le nb $z\in\mathcal D$ choisi, mais par le nb complexe $\tau=z_2/z_1$ lui-même (symétrique de $z$ p.r. à l'axe imaginaire), on trouverait une formule de transformation d'opération de $\mathrm{Sl}(2,\mathbb Z)$ sur $\mathcal D$ qui n'est \uline{pas} la formule habituelle.

\margin{Donc cela montre que la condition $z\in\mathcal D^+$ signifie que $(z_1,z_2)$ est, non une base directe, mais une base \uline{inverse} de $\unsure{t_E}$. Si on \deleted{voulait}\inserted{veut} faire correspondre à $z\in\mathcal D^+$ une structure de Torelli (directe), il y a lieu de changer la convention (23) en $L_z=\mathbb C(1,z)$ qui donne une \deleted{\unsure{orientation}}\inserted{structure complexe sur $\mathbb R^2$} compatible avec l'orientation canonique (car $L_z\in\Sigma_{\mathbb R^2,\omega_0}^{-}$, cf \unsure{calcul ci-dessus}). On trouve
\[
\uline{z=-(z_2/z_1)^{-1}=-\frac{z_1}{z_2}}.
\]
C'est un point \unsure{auquel il|où il} faudrait vérifier que le foncteur en question est bien analytiquement naturel\ldots}

On trouve ici le

\uline{Théorème} $\mathcal D$ \inserted{$=\{z\in\mathbb C\mid\Im z>0\}$,} muni de la famille \inserted{analytique} de courbes elliptiques
\[
E_z=\mathbb C/(\mathbb Z\cdot1+\mathbb Z\cdot(-\overline z))
\]
est un revêtement universel de $M_{1,1}^{\mathrm{an}}$, correspondant à la classification des courbes elliptiques avec rigidification de Torelli; le groupe des $M_{1,1}^{\mathrm{an}}$-automorphismes de $\mathcal D$ (opérant à droite) s'identifie au groupe $\mathrm{Sl}(2,\mathbb Z)$, opérant sur $\mathcal D$ par
\[
gz=\frac{az+b}{cz+d}\qquad\text{pour }\quad
g=\begin{pmatrix}a&b\\c&d\end{pmatrix}\in\mathrm{Sl}(2,\mathbb Z),
\]
qui correspond ainsi à l'opération à g. naturelle
